\documentclass[11pt]{ctexart}

% 中文版报告：使用 ctexart 编译，建议 XeLaTeX。
% ============================================================
% setup.tex
% Packages, formatting, commands, and reusable LaTeX structures
% ============================================================

% -------------------------
% Engine-safe font setup
% -------------------------
\usepackage{iftex}
\ifPDFTeX
    \usepackage[T1]{fontenc}
    \usepackage[utf8]{inputenc}
\else
    % XeLaTeX/LuaLaTeX use UTF-8 natively.
    % We intentionally avoid explicit font selection for maximum Overleaf compatibility.
\fi

% -------------------------
% Page Layout
% -------------------------
\usepackage[margin=1in]{geometry}
\usepackage{setspace}
\setstretch{1.08}
\setlength{\parskip}{0.35em}
\setlength{\parindent}{0pt}
\emergencystretch=3em
\sloppy

% -------------------------
% Math
% -------------------------
\usepackage{amsmath}
\usepackage{amssymb}
\usepackage{mathtools}

% -------------------------
% Figures and Tables
% -------------------------
\usepackage{graphicx}
\usepackage{booktabs}
\usepackage{multirow}
\usepackage{array}
\usepackage{caption}
\usepackage{subcaption}
\usepackage{float}
\usepackage{placeins}

% -------------------------
% Lists
% -------------------------
\usepackage{enumitem}
\setlist[itemize]{leftmargin=1.5em}
\setlist[enumerate]{leftmargin=1.5em}

% -------------------------
% Algorithms
% -------------------------
\usepackage{algorithm}
\usepackage{algpseudocode}

% -------------------------
% Colors and Hyperlinks
% -------------------------
\usepackage{xcolor}
\usepackage[
    colorlinks=true,
    linkcolor=blue,
    citecolor=blue,
    urlcolor=blue
]{hyperref}

% -------------------------
% Code Blocks
% -------------------------
\usepackage{listings}

\definecolor{codegray}{RGB}{245,245,245}
\definecolor{commentgray}{RGB}{110,110,110}
\definecolor{keywordblue}{RGB}{0,70,160}
\definecolor{stringteal}{RGB}{0,120,120}

\lstdefinestyle{codestyle}{
    backgroundcolor=\color{codegray},
    basicstyle=\ttfamily\small,
    breaklines=true,
    frame=single,
    numbers=left,
    numberstyle=\tiny\color{commentgray},
    keywordstyle=\color{keywordblue},
    commentstyle=\color{commentgray},
    stringstyle=\color{stringteal},
    showstringspaces=false,
    tabsize=4,
    columns=fullflexible,
    keepspaces=true,
    literate={├}{{|}}1 {└}{{|}}1 {─}{{-}}1 {│}{{|}}1
}

\lstset{style=codestyle}

% -------------------------
% Convenient Commands
% -------------------------
\newcommand{\methodname}{Offline OPD}
\newcommand{\taskname}{Multi-turn Troubleshooting Agent}
\newcommand{\supportopd}{Support-aware Offline OPD}

% Avoid bold-small-caps font warnings by keeping method names plain.
\newcommand{\sft}{SFT}
\newcommand{\onlinerl}{Online RL}
\newcommand{\offlineopd}{Offline OPD}
\newcommand{\onlineopd}{Online OPD}

% Safe code command: raw underscores are allowed inside \code{...}.
% Example: \code{offline_opd.py} compiles correctly.
\newcommand{\code}[1]{\texttt{\detokenize{#1}}}

% -------------------------
% Placeholder Figure Command
% This avoids compilation failure when a figure file is missing.
% The \detokenize call prevents underscores in filenames from breaking LaTeX.
% -------------------------
\newcommand{\placeholderfigure}[3]{%
\begin{figure}[H]
\centering
\IfFileExists{#1}{%
    \includegraphics[width=#2\linewidth]{#1}%
}{%
    \fbox{%
        \begin{minipage}[c][0.25\textheight][c]{#2\linewidth}
        \centering
        Placeholder for figure: \texttt{\detokenize{#1}}
        \end{minipage}%
    }%
}%
\caption{#3}
\end{figure}%
}

% -------------------------
% Table Column Helpers
% -------------------------
\newcolumntype{L}[1]{>{\raggedright\arraybackslash}p{#1}}
\newcolumntype{C}[1]{>{\centering\arraybackslash}p{#1}}

% -------------------------
% Document Metadata Helper
% -------------------------
\newcommand{\reporttitle}{%
    \textbf{Offline OPD under Agentic Shift} \\
    \large A Lightweight Multi-turn Agentic Toy Task Study%
}

\newcommand{\reportauthor}{%
    Runze Lv \\
    \texttt{runzelvlc@gmail.com}%
}


\title{\textbf{Offline OPD under Agentic Shift}\\
\large 离线 OPD 在多轮智能体分布偏移下的假设边界实验}
\author{Runze Lv\\\texttt{runzelvlc@gmail.com}}
\date{2026 年 6 月}

\begin{document}

\maketitle

\begin{abstract}
本文选择在线考核题目二：\textbf{Offline OPD under Agentic Shift}。
围绕 Lightning OPD 的离线化假设，本文构造了一个 CPU-only 的多轮故障排查 toy task。
在该任务中，智能体需要通过工具式动作逐步诊断隐藏故障，并提交最终根因；错误的早期修复动作会污染日志或误配置系统组件，从而改变后续状态分布。
本文实现并比较了 \sft{}、\onlinerl{}、\offlineopd{}、\onlineopd{}，并额外实现了一个 support-aware \offlineopd{} patch。
实验结果显示：在正常分布下，\sft{} 与 \offlineopd{} 都能达到很高成功率；但在注入早期错误动作的 agentic-shift 压力测试中，off-support state ratio 明显上升，成功率下降，从而暴露了离线 OPD 对状态覆盖的依赖。
\end{abstract}

\section{题目选择与问题理解}

本文选择的题目是第二题：\textbf{Offline OPD under Agentic Shift｜离线 OPD 的假设边界}。
题目要求阅读 Lightning OPD 及对比论文，理解 \sft{}、\onlinerl{}、\onlineopd{}、\offlineopd{} 的区别，并设计一个轻量级多轮 agentic toy task。
该任务需要满足以下条件：

\begin{itemize}
    \item 多轮决策；
    \item 有工具式动作；
    \item 有错误分支；
    \item 有最终答案奖励；
    \item 早期错误会影响后续状态。
\end{itemize}

本文代码实现位于项目根目录的 \code{opd_agentic_shift/} 下。
中文报告源码位于 \code{latex_repo/} 下，提交用的汇总结果位于项目根目录的 \code{results/} 下。

\section{论文背景与方法差异}

\subsection{Lightning OPD 的核心假设}

Lightning OPD 的目标是减少在线 teacher query 的成本。
它希望利用已经由初始策略生成的 rollout 数据，在这些状态上预先计算 teacher 的动作分布或 log-probability，然后离线训练学生策略。
如果 teacher consistency 假设成立，即同一状态上的 teacher label 可以被稳定复用，那么 \offlineopd{} 就可以避免训练时反复调用 teacher。

但是这个假设在多轮 agentic 任务中会受到挑战。
原因是策略的一个早期动作可能改变环境状态，导致后续状态不再落在原来的离线数据支持集内。
此时，离线缓存的 teacher label 只能覆盖旧分布，无法回答新分支上的 teacher 应该怎么做。

\subsection{与对比论文的区别}

对比论文 \emph{Revealing the Power of Post-Training for Small Language Models via Knowledge Distillation} 也使用了 SFT 与离线蒸馏思想。
它更偏向小模型 post-training 的实用流程：先进行 curriculum SFT，再通过 teacher-student knowledge distillation 提升模型能力。
Lightning OPD 更强调从 on-policy distillation 的角度分析 teacher consistency 与离线化假设。

简言之：

\begin{itemize}
    \item 对比论文更关注小模型蒸馏训练 recipe；
    \item Lightning OPD 更关注如何把 OPD 从在线 teacher query 转为离线训练；
    \item 本文实验则进一步考察：当任务变成多轮工具调用，并且早期动作会改变后续状态时，离线 OPD 的假设边界在哪里。
\end{itemize}

\section{Toy Task：多轮故障排查智能体}

\subsection{任务定义}

本文设计了一个符号化的多轮故障排查环境。
每个 episode 开始时，环境从 5 类隐藏故障中采样一个真实根因：

\begin{table}[H]
\centering
\caption{隐藏故障类型}
\begin{tabular}{ll}
\toprule
编号 & 隐藏故障 \\
\midrule
1 & \code{network_dns_error} \\
2 & \code{disk_full} \\
3 & \code{database_connection_error} \\
4 & \code{permission_denied} \\
5 & \code{service_crash} \\
\bottomrule
\end{tabular}
\end{table}

智能体不能直接观察真实故障，只能观察工具调用返回的状态，例如日志状态、网络状态、磁盘状态、数据库状态和权限状态。
智能体最终需要选择一个 \code{submit_*} 动作提交根因。

\subsection{动作空间}

动作空间分为三类：诊断工具、修复动作和最终答案。

\begin{table}[H]
\centering
\caption{故障排查环境的动作空间}
\begin{tabular}{L{0.25\linewidth} L{0.65\linewidth}}
\toprule
类别 & 动作 \\
\midrule
诊断工具
&
\code{check_service_log},
\code{check_network},
\code{check_disk},
\code{check_db},
\code{check_permission}
\\
修复动作
&
\code{restart_service},
\code{clean_disk},
\code{fix_dns},
\code{fix_db_config},
\code{fix_permission}
\\
最终答案
&
\code{submit_network_dns_error},
\code{submit_disk_full},
\code{submit_database_connection_error},
\code{submit_permission_denied},
\code{submit_service_crash}
\\
\bottomrule
\end{tabular}
\end{table}

\subsection{早期错误造成的状态偏移}

为了模拟 agentic shift，错误的修复动作会产生持久副作用：

\begin{itemize}
    \item 过早执行 \code{restart_service} 会污染日志；
    \item 非 DNS 故障时执行 \code{fix_dns} 会造成网络误配置；
    \item 非数据库故障时执行 \code{fix_db_config} 会造成数据库误配置；
    \item 非权限故障时执行 \code{fix_permission} 会造成权限相关误配置；
    \item 某些错误清理动作会让后续日志变得不可靠。
\end{itemize}

因此，智能体的后续观测不仅取决于隐藏故障，也取决于此前动作。
这使得任务具备多轮 agentic 系统中常见的错误累积问题。

\subsection{奖励函数}

奖励函数以最终诊断为主：

\begin{equation}
r_t =
\begin{cases}
+1.0, & \text{提交正确根因}, \\
-1.0, & \text{提交错误根因}, \\
-0.05, & \text{诊断工具调用}, \\
-0.20, & \text{错误修复动作}, \\
-0.50, & \text{超时未完成}.
\end{cases}
\end{equation}

最大轨迹长度设为 $T=6$。

\section{代码框架}

项目主体结构如下：

\begin{lstlisting}
Offline_OPD/
├── run_all.py
├── README.md
├── REPORT.md
├── results/
│   ├── eval/
│   ├── eval_shift/
│   └── tables/
├── opd_agentic_shift/
│   ├── envs/
│   │   ├── troubleshooting_env.py
│   │   └── encoder.py
│   ├── teachers/
│   │   └── rule_teacher.py
│   ├── policies/
│   │   └── mlp_policy.py
│   ├── data/
│   │   ├── generate_expert_data.py
│   │   ├── collect_rollouts.py
│   │   └── build_offline_opd_data.py
│   ├── algos/
│   │   ├── sft.py
│   │   ├── offline_opd.py
│   │   ├── online_rl.py
│   │   └── online_opd.py
│   ├── eval/
│   │   ├── evaluate.py
│   │   └── case_analysis.py
│   └── utils/
│       └── io.py
└── latex_repo/
    ├── main_cn.tex
    ├── setup.tex
    └── results/
\end{lstlisting}

\begin{table}[H]
\centering
\caption{主要代码模块}
\begin{tabular}{L{0.32\linewidth} L{0.20\linewidth} L{0.38\linewidth}}
\toprule
文件 & 角色 & 说明 \\
\midrule
\code{troubleshooting_env.py} & 环境 & 多轮故障排查环境，包含错误分支和副作用 \\
\code{encoder.py} & 状态编码 & 将符号观测转为 MLP 输入向量 \\
\code{rule_teacher.py} & Teacher & 可读取隐藏故障的规则 oracle \\
\code{mlp_policy.py} & Policy & 小型 MLP 学生策略 \\
\code{generate_expert_data.py} & 数据生成 & 用 teacher 生成专家轨迹 \\
\code{collect_rollouts.py} & Rollout 收集 & 用 SFT policy 收集 student states \\
\code{build_offline_opd_data.py} & 离线标注 & 在 student states 上预计算 teacher distribution \\
\code{sft.py} & 训练 & 行为克隆 baseline \\
\code{online_rl.py} & 训练 & REINFORCE baseline \\
\code{offline_opd.py} & 训练 & 普通离线 OPD 与 support-aware patch \\
\code{online_opd.py} & 训练 & 在线 teacher query 的上界方法 \\
\code{evaluate.py} & 评估 & success、reward、length、off-support ratio \\
\bottomrule
\end{tabular}
\end{table}

\section{方法实现}

\subsection{SFT}

\sft{} 使用 rule teacher 生成的专家轨迹训练学生策略：

\begin{equation}
\mathcal{L}_{\text{SFT}}(\theta)
=
-\mathbb{E}_{(s,a^*) \sim \mathcal{D}_{\text{expert}}}
\log \pi_\theta(a^* \mid s).
\end{equation}

该方法稳定，但只见过专家访问的状态。
一旦学生策略在测试时进入错误分支，就可能遇到 expert data 不覆盖的状态。

\subsection{Online RL}

\onlinerl{} 使用 REINFORCE 直接优化环境 reward：

\begin{equation}
\mathcal{L}_{\text{RL}}(\theta)
=
-\mathbb{E}_{\tau \sim \pi_\theta}
\sum_t \log \pi_\theta(a_t \mid s_t) R_t.
\end{equation}

它可以访问自身策略诱导出的状态，但在稀疏奖励下训练更不稳定。

\subsection{Offline OPD}

\offlineopd{} 首先用 SFT policy 收集 rollout，然后在这些状态上预先计算 teacher 分布：

\begin{equation}
\mathcal{D}_{\text{OPD}}
=
\{(s_t, q_T(\cdot \mid s_t))\}_{t=1}^{N}.
\end{equation}

学生策略通过匹配 teacher 分布训练：

\begin{equation}
\mathcal{L}_{\text{Offline OPD}}(\theta)
=
\mathbb{E}_{s \sim \mathcal{D}_{\text{OPD}}}
\left[
\mathrm{KL}
\left(
q_T(\cdot \mid s)
\parallel
\pi_\theta(\cdot \mid s)
\right)
\right].
\end{equation}

该方法比专家轨迹 SFT 更接近学生策略分布，但仍受限于固定 offline rollout 的覆盖范围。

\subsection{Online OPD}

\onlineopd{} 在训练时让当前学生策略 rollout，并实时查询 teacher：

\begin{equation}
s_t \sim \pi_\theta,
\quad
q_T(\cdot \mid s_t) = \mathrm{Teacher}(s_t).
\end{equation}

它可以作为一个上界，因为 teacher label 始终来自当前学生分布。
但实际大模型场景中，在线 teacher query 成本较高。

\section{提出的 Patch：Support-aware Offline OPD}

本文针对 offline state coverage 不足提出一个轻量 patch：\supportopd{}。
其核心思想是根据状态在离线 OPD 数据集中出现的频率，为每个样本赋权：

\begin{equation}
w(s)
=
\min
\left(
1,
\frac{\mathrm{count}(s)}{\tau}
\right)^{1/2},
\end{equation}

其中 $\tau$ 是 support threshold。
最终 loss 为：

\begin{equation}
\mathcal{L}_{\text{support}}
=
\mathbb{E}_{s \sim \mathcal{D}_{\text{OPD}}}
\left[
w(s)
\cdot
\mathrm{KL}
\left(
q_T(\cdot \mid s)
\parallel
\pi_\theta(\cdot \mid s)
\right)
\right]
+
\lambda
\mathrm{KL}
\left(
\pi_{\text{SFT}}(\cdot \mid s)
\parallel
\pi_\theta(\cdot \mid s)
\right).
\end{equation}

直觉上，低覆盖状态上的 teacher label 可能不足以代表真实测试分布，因此训练时更保守；
同时 KL anchor 防止策略过度偏离 SFT 初始化。

\section{实验设置}

一键运行命令如下：

\begin{lstlisting}
python run_all.py --profile quick
\end{lstlisting}

快速实验配置如下：

\begin{table}[H]
\centering
\caption{Quick profile 实验设置}
\begin{tabular}{lr}
\toprule
设置项 & 数值 \\
\midrule
expert episodes & 300 \\
SFT rollout episodes & 300 \\
SFT epochs & 12 \\
Offline OPD epochs & 15 \\
Online RL episodes & 600 \\
Online OPD episodes & 300 \\
evaluation episodes & 200 \\
\bottomrule
\end{tabular}
\end{table}

评估指标包括：

\begin{itemize}
    \item \code{success_rate}：最终诊断正确率；
    \item \code{avg_reward}：平均 episode reward；
    \item \code{avg_length}：平均轨迹长度；
    \item \code{off_support_vs_expert}：相对于 expert states 的 off-support 比例；
    \item \code{off_support_vs_offline_opd}：相对于 offline OPD 数据集的 off-support 比例。
\end{itemize}

除正常评估外，本文还加入了 agentic-shift 压力测试。
在该测试中，每个 episode 开始后先注入一次错误 repair action，再让 policy 接管。
这用于模拟多轮智能体中常见的早期错误状态偏移。

\section{实验结果}

\subsection{正常分布评估}

\begin{table}[H]
\centering
\caption{In-distribution evaluation results}
\begin{tabular}{lrrrrr}
\toprule
方法 & 成功率 & 平均奖励 & 平均长度 & Off-support/expert & Off-support/OPD \\
\midrule
\sft{} & 1.000 & 0.861 & 2.785 & 0.000 & 0.000 \\
\onlinerl{} & 0.755 & 0.410 & 2.000 & 0.000 & 0.000 \\
\offlineopd{} & 1.000 & 0.861 & 2.785 & 0.000 & 0.000 \\
\supportopd{} & 1.000 & 0.861 & 2.785 & 0.000 & 0.000 \\
\onlineopd{} & 1.000 & 0.861 & 2.785 & 0.000 & 0.000 \\
\bottomrule
\end{tabular}
\end{table}

正常分布下，专家路径覆盖了主要状态，\sft{}、\offlineopd{}、\supportopd{} 和 \onlineopd{} 都能达到 100\% 成功率。
\onlinerl{} 因为奖励稀疏且训练 episode 较少，表现较弱。
这一结果说明：当状态分布稳定、离线数据覆盖充分时，offline OPD 的离线化假设基本成立。

\subsection{Agentic-shift 压力测试}

\begin{table}[H]
\centering
\caption{Injected early wrong repair stress results}
\begin{tabular}{lrrrrr}
\toprule
方法 & 成功率 & 平均奖励 & 平均长度 & Off-support/expert & Off-support/OPD \\
\midrule
\sft{} & 0.790 & 0.188 & 2.850 & 1.000 & 0.919 \\
\onlinerl{} & 0.580 & -0.183 & 1.865 & 1.000 & 0.949 \\
\offlineopd{} & 0.870 & 0.357 & 3.055 & 1.000 & 0.925 \\
\supportopd{} & 0.870 & 0.357 & 3.055 & 1.000 & 0.925 \\
\onlineopd{} & 0.730 & 0.177 & 3.350 & 1.000 & 0.931 \\
\bottomrule
\end{tabular}
\end{table}

压力测试中，所有方法的 \code{off_support_vs_expert} 都达到 1.000，说明注入的错误 repair action 确实将 policy 带入了专家轨迹没有覆盖的状态。
相对于 offline OPD rollout，off-support ratio 也在 0.919 到 0.949 之间，说明学生遇到的大量状态并不在离线 OPD 数据支持内。

这一结果直接说明：多轮 agentic 任务中，早期错误会显著破坏 offline coverage。
即使 offline OPD 在正常分布下表现很好，在 shifted states 上也会受到固定离线数据的限制。

\section{Case 分析}

\subsection{成功 case}

在正常分布中，\sft{} 的一个成功案例为：

\begin{lstlisting}
fault = service_crash
actions = [
  "check_service_log",
  "submit_service_crash"
]
reward = 0.90
\end{lstlisting}

该案例中，日志工具直接返回了 service crash 线索，因此策略能够快速提交正确答案。

\subsection{失败 case}

\onlinerl{} 在正常分布中的一个失败案例为：

\begin{lstlisting}
fault = disk_full
actions = [
  "check_service_log",
  "submit_network_dns_error"
]
reward = -1.10
\end{lstlisting}

该策略过早提交了错误根因，体现出 sparse reward 下 online RL 的不稳定性。

\subsection{Shifted failure case}

在 agentic-shift 压力测试中，\offlineopd{} 的一个失败案例为：

\begin{lstlisting}
fault = database_connection_error
injected_action = clean_disk
policy_actions = [
  "check_service_log",
  "check_network",
  "check_permission",
  "submit_network_dns_error"
]
reward = -1.45
\end{lstlisting}

错误的早期动作污染了后续观测，使策略进入离线数据低覆盖分支。
随后策略检查了多个无关工具，并提交了错误根因。
这正是 offline OPD 在 agentic shift 下可能失败的典型路径。

\section{结论}

本文实现了一个轻量级、可复现的多轮 agentic toy task，用于考察 \offlineopd{} 的假设边界。
实验结论如下：

\begin{itemize}
    \item 当测试状态分布与 SFT rollout 分布接近时，\offlineopd{} 可以稳定工作；
    \item 当早期错误动作改变后续状态时，off-support ratio 显著上升，离线 OPD 的固定数据覆盖成为瓶颈；
    \item \onlinerl{} 能访问自身诱导状态，但在小样本和稀疏奖励下不稳定；
    \item \onlineopd{} 理论上可以标注当前学生分布，但需要在线 teacher query；
    \item support-aware loss 是一个低成本 patch，可以表达对低覆盖状态的保守性，但如果离线数据本身覆盖不足，它不能完全替代 DAgger-style refresh 或在线 teacher query。
\end{itemize}

因此，Lightning OPD 的离线化假设在单轮或短输出任务中更容易成立；
而在多轮工具调用任务中，需要额外考虑 state coverage、错误分支累积和 teacher query refresh 策略。

\section{复现命令与结果文件}

安装依赖：

\begin{lstlisting}
pip install -r requirements.txt
\end{lstlisting}

快速复现实验：

\begin{lstlisting}
python run_all.py --profile quick
\end{lstlisting}

更大规模复现实验：

\begin{lstlisting}
python run_all.py --profile full
\end{lstlisting}

本文使用的汇总结果已经归档到：

\begin{lstlisting}
results/
├── eval/
│   └── summary.json
├── eval_shift/
│   └── summary.json
└── tables/
    └── summary_tables_cn.csv
\end{lstlisting}

完整训练数据、checkpoint 和逐 episode 详情保留在项目根目录的 \code{runs/} 下，可通过上述命令重新生成。

\begin{thebibliography}{9}

\bibitem{lightningopd}
Lightning OPD.
\newblock \url{https://arxiv.org/pdf/2604.13010}

\bibitem{posttrainingkd}
Revealing the Power of Post-Training for Small Language Models via Knowledge Distillation.
\newblock \url{https://arxiv.org/pdf/2509.26497}

\end{thebibliography}

\end{document}
